% !TeX root = ../../Book3.tex
% 纲要标题：### 24.1 结构层
% 纲要位置：计划纲要.md，第 2736 行。

\section{结构层}
\label{b3:sec:2401}


% 纲要内容（仅作写作计划，尚非教材正文）：
%
% * 基本开限制的方向及非单射情形；环值、集合值和群值层采用同一黏合公理
% * 空覆盖、芽与剩余域值的区别；由局部分式族建立结构层，再证明基本开截面与茎公式
% * 局部同态保持指定点的消失条件；仿射态射对应、Noether 概形与拟紧性的条件
%
% ---
%

前面已经把环的素理想组织成拓扑空间，但拓扑本身无法区分约化点与加厚点。本节把各个基本开集上的局部化同时保留下来，并说明这些环如何限制、黏合。这种记录邻域中代数信息的方法，就是结构层。

\subsection{基本开集上的局部化}
\label{b3:subsec:2401:01}

在交换环 \(A\) 的素谱上，\cref{b3:thm:spec-localization}把 \(D(f)\) 识别为 \(\operatorname{Spec}A_f\)。把开集缩小以后，原来尚不可逆的元素可能变得可逆，所以截面应从较大开集限制到较小开集。具体地，若 \(D(g)\subseteq D(f)\)，则 \(f\) 在 \(A_g\) 中可逆：包含关系给 \(g^N=af\)，从而 \(a/g^N\) 是其逆。因此有唯一限制同态 \(A_f\rightarrow A_g\)，且这些同态与复合相容。

\ProseParagraphBreak
这种限制同态不一定单射，因为缩小开集也可能使原来的元素消失。例如，对域 \(k\)，令 \(A=k\times k\)、\(e=(1,0)\)，则 \(A\to A_e\cong k\) 是第一分量投影，它把非零的 \((0,1)\) 送到零。这里记录的是限制同态，不能把不同开集上的环一律放成包含关系。从相容的局部分式恢复整体元素所需的代数步骤，已由\cref{b3:lem:localization-sheaf-gluing}证明。

\ProseParagraphBreak
对 \(A_f\) 使用该引理，可以处理 \(D(f)\) 的基本开覆盖；任意这样的覆盖由\cref{b3:prop:spec-quasicompact}取有限子覆盖，未选中的开集再由唯一性核对。因此局部化既能限制，也能在交集相容时黏合。这正是下一小节层公理的代数来源。

\subsection{仿射概形的结构层}
\label{b3:subsec:2401:02}

先写本章所需的环值情形。随后将值换成集合或交换群时，限制与唯一黏合仍采用同一套公理，只是保留的代数运算不同。

\begin{definition}\label{b3:def:sheaf-rings}
设 \(X\) 为拓扑空间。对每个开集 \(U\subseteq X\) 给定交换环 \(\mathcal F(U)\)，对 \(V\subseteq U\) 给定保持单位的限制同态 \(\rho_{UV}:\mathcal F(U)\rightarrow\mathcal F(V)\)。若 \(\rho_{UU}=\operatorname{id}\)，且对 \(W\subseteq V\subseteq U\) 有 \(\rho_{VW}\rho_{UV}=\rho_{UW}\)，则称这些数据为环的\term[yuceng]{预层}[presheaf]。若再对任意开覆盖 \(U=\bigcup_iU_i\)，在两两交集上限制相同的 \(s_i\in\mathcal F(U_i)\) 都来自唯一的 \(s\in\mathcal F(U)\)，则称其为环的\term[ceng]{层}[sheaf]。空覆盖约定使 \(\mathcal F(\varnothing)\) 为零环。元素 \(s\in\mathcal F(U)\) 称为 \(U\) 上的截面，也记 \(\Gamma(U,\mathcal F)\)。
\end{definition}
\symidx{\(\Gamma(U,\mathcal F)\)：层 \(\mathcal F\) 在开集 \(U\) 上的截面}

\begin{definition}\label{b3:def:sheaf-sets-groups}
设 \(X\) 为拓扑空间。对每个开集 \(U\subseteq X\) 给定集合 \(\mathcal F(U)\)，对每个包含关系 \(V\subseteq U\) 给定限制映射 \(\rho_{UV}:\mathcal F(U)\rightarrow\mathcal F(V)\)，并记 \(s|_V=\rho_{UV}(s)\)。若 \(\rho_{UU}=\operatorname{id}\)，且对 \(W\subseteq V\subseteq U\) 有 \(\rho_{VW}\rho_{UV}=\rho_{UW}\)，则称这些数据为集合值预层。若再对任意开覆盖 \(U=\bigcup_iU_i\)，满足
\[
s_i|_{U_i\cap U_j}=s_j|_{U_i\cap U_j}
\quad\text{对所有}\;i,j\;\text{成立}
\]
的截面族 \(s_i\in\mathcal F(U_i)\) 都来自唯一的 \(s\in\mathcal F(U)\)，则称其为集合值层。若各 \(\mathcal F(U)\) 为交换群，且限制映射为群同态，则在相同公理下分别称为交换群值预层、交换群值层。对层，空集的空覆盖要求 \(\mathcal F(\varnothing)\) 恰有一个元素；在交换群值情形，这就是零群。
\end{definition}

环层忘掉乘法以后是交换群值层，再忘掉加法以后是集合值层，因为限制与唯一黏合的条件都保留下来。上一小节的 \(A_f\) 因而同时给出三种取值的局部模型；对 \(A\) 模 \(M\)，\cref{b3:lem:localization-sheaf-gluing}中的 \(M_f\) 也给出基本开集上的交换群值黏合模型。空基本开集 \(D(0)\) 上分别取零环、零群或单点集，与空覆盖的要求一致。

\ProseParagraphBreak

集合值层的态射是一组与限制相容的集合映射；对交换群值层、环层，则分别要求群同态、保持单位的环同态。这三种情形都可以谈截面、芽和茎：若两个截面分别定义在含 \(x\) 的开邻域上，并在某个更小的含 \(x\) 邻域上限制相等，就定义相同的芽。所有这样的芽组成茎 \(\mathcal F_x=\varinjlim_{x\in U}\mathcal F(U)\)。在交换群值或环值情形，先把代表截面限制到共同邻域，再作运算；限制保留运算，保证所得结果与代表截面的选择无关，因而得到茎上的交换群或环结构。芽只保留任意充分小邻域内的信息，所以不是把截面简单地“代入点”所得的数。
\symidx{\(\mathcal F_x\)：层 \(\mathcal F\) 在点 \(x\) 处的茎}

\ProseParagraphBreak
一般开集未必形如一个 \(D(f)\)，因而没有预先给定的单一局部化环。结构层先要求每一点附近都能用同一个分式表示，再把这些局部表示黏合到任意开集上；下面的基本开截面定理将说明，这个构造在 \(D(f)\) 上恰好返回熟悉的 \(A_f\)。

\begin{definition}\label{b3:def:affine-structure-sheaf}
设 \(A\) 为交换环，\(X=\operatorname{Spec}A\)。对开集 \(U\subseteq X\)，令 \(\mathcal O_X(U)\) 由所有族 \((s_{\mathfrak p})_{\mathfrak p\in U}\)、\(s_{\mathfrak p}\in A_{\mathfrak p}\) 组成，并要求每个点都有开邻域 \(V\subseteq U\) 及 \(a,b\in A\)，使所有 \(\mathfrak q\in V\) 都满足 \(b\notin\mathfrak q\) 及 \(s_{\mathfrak q}=a/b\)。配以逐点运算和通常限制，所得层称为素谱的\term[jiegou-ceng]{结构层}[structure sheaf]。
\end{definition}
\symidx{\(\mathcal O_X\)：仿射概形的结构层}

限制保持局部表示；相容的族可按点拼合，拼合后仍在原邻域上使用原分式，因而确实满足层公理。这里的值在局部环 \(A_{\mathfrak p}\) 中，而非仅在剩余域 \(\kappa(\mathfrak p)\) 中，这使结构层仍能记录幂零元。

\begin{theorem}\label{b3:thm:affine-sheaf-sections}
设 \(A\) 为交换环，\(X=\operatorname{Spec}A\)，\(f\in A\)，\(\mathfrak p\in X\)。自然映射给出
\[
\Gamma(D(f),\mathcal O_X)\cong A_f,\qquad
\mathcal O_{X,\mathfrak p}\cong A_{\mathfrak p}.
\]
特别地，\(\Gamma(X,\mathcal O_X)\cong A\)；这些同构与限制和局部化相容。
\end{theorem}
\begin{proof}
先识别茎与局部环，再证明基本开截面映射单射，最后用有限子覆盖上的分式黏合证明满射。三个步骤分别使用局部化的相等判据、局部零检测与清分母黏合。

先看茎。分式 \(a/b\)、\(b\notin\mathfrak p\)，在 \(D(b)\) 上给出截面，所以每个局部环元素都来自芽。若两个分式在 \(A_{\mathfrak p}\) 中相等，存在 \(c\notin\mathfrak p\) 清除它们差的分子，故它们在 \(D(b_1b_2c)\) 上相等。这证明芽与局部环一一对应。

自然映射 \(A_f\rightarrow\Gamma(D(f),\mathcal O_X)\) 是单射。事实上，若一个分式在每个 \(A_{\mathfrak p}\)、\(\mathfrak p\in D(f)\)，中为零，按分式零判据可用一组更小的基本开集使它为零；取有限子覆盖后，由\cref{b3:lem:localization-sheaf-gluing}的单射部分，在 \(A_f\) 中也为零。

给定 \(D(f)\) 上截面，将其局部表示邻域缩为基本开集 \(D(g_i)\)。若局部表达的分母为 \(b_i\)，则 \(D(g_i)\subseteq D(b_i)\) 保证 \(b_i\) 在 \(A_{g_i}\) 可逆，故该局部截面确实来自 \(A_{g_i}\)。有限子覆盖上的这些元素在交集一致；刚证的单射性把截面相等变成局部化中的相等。对 \(A_f\) 应用黏合引理，便得到唯一的 \(A_f\) 元素。所有映射都由同一分式给出，故自动与限制相容。
\end{proof}

例如，对域 \(k\)，在 \(\operatorname{Spec}k[x]\) 的原点 \((x)\) 处，分式 \(1/(1-x)\) 定义在含原点的 \(D(1-x)\) 上，给出茎 \(k[x]_{(x)}\) 中的一个芽。它不是全局多项式；但模去局部环的极大理想后，它的值为 \(1\in k\)。常数 \(1\) 与这个分式具有相同的剩余域值，却给出不同的芽，差为非零的 \(x/(1-x)\)。

\ProseParagraphBreak
\(\operatorname{Spec}k\) 与 \(\operatorname{Spec}k[\varepsilon]/(\varepsilon^2)\) 都只有一个点，结构层的全局截面却分别为 \(k\) 与 \(k[\varepsilon]/(\varepsilon^2)\)。因此相同的拓扑空间可以承载不同的代数结构。结构层的基本性质可对照 Hartshorne 的《Algebraic Geometry》\cite[Chapter II, Proposition 2.2, pp. 71--72]{Hartshorne1977}。

\subsection{局部环空间与态射}
\label{b3:subsec:2401:03}

\begin{definition}\label{b3:def:locally-ringed-space}
设 \(X\) 为拓扑空间，\(\mathcal O_X\) 为其上的交换环层。若每个茎 \(\mathcal O_{X,x}\) 都是局部环，则称 \((X,\mathcal O_X)\) 为\term[jubu-huan-kongjian]{局部环空间}[locally ringed space]。从 \((X,\mathcal O_X)\) 到 \((Y,\mathcal O_Y)\) 的态射由连续映射 \(f:X\rightarrow Y\) 及与限制相容的同态
\[
f^\#_V:\mathcal O_Y(V)\longrightarrow\mathcal O_X(f^{-1}V)
\]
组成，且要求每个诱导茎同态 \(f_x^\#:\mathcal O_{Y,f(x)}\rightarrow\mathcal O_{X,x}\) 为局部同态。若两局部环的极大理想分别记为 \(\mathfrak m_{f(x)}\)、\(\mathfrak m_x\)，这要求
\[
(f_x^\#)^{-1}(\mathfrak m_x)=\mathfrak m_{f(x)}.
\]
\end{definition}

一个芽在目标点的剩余域中为零，当且仅当其拉回在源点的剩余域中为零；在目标点可逆的芽，拉回后仍可逆。因此局部同态不仅给出邻域函数的拉回，也使这一拉回尊重所指定的点。

\begin{definition}\label{b3:def:scheme}
设 \((X,\mathcal O_X)\) 为局部环空间。若存在交换环 \(A\) 及局部环空间同构
\[
(X,\mathcal O_X)\cong(\operatorname{Spec}A,\mathcal O_{\operatorname{Spec}A}),
\]
则称 \(X\) 为\term[fangshe-gaixing]{仿射概形}[affine scheme]。若 \(X\) 的每点都有开邻域 \(U\)，使 \((U,\mathcal O_X|_U)\) 为仿射概形，则称 \(X\) 为\term[gaixing]{概形}[scheme]。概形态射就是局部环空间态射。
\end{definition}

在这个定义中，\(\mathcal O_X|_U\) 表示只在 \(U\) 的开集上取截面。由\cref{b3:thm:affine-sheaf-sections}，仿射概形的基本开集仍为仿射概形；一般开子概形则不必仿射。

\begin{definition}\label{b3:def:noetherian-scheme}
设 \(X\) 为概形。若 \(X\) 有由 Noether 环的素谱组成的仿射开覆盖，则称 \(X\) 为局部 Noether 概形；若这样的覆盖还可以取为有限覆盖，则称 \(X\) 为 Noether 概形。后一个条件等价于局部 Noether 且底空间拟紧。
\end{definition}

若 \(X\) 局部 Noether 且拟紧，从定义中的 Noether 仿射开覆盖抽取有限子覆盖，即得到 Noether 概形。反之，每个仿射谱拟紧，有限个这样的开集覆盖 \(X\) 时，任意开覆盖都可在每个仿射块上取有限子覆盖，再合成 \(X\) 的有限子覆盖，所以底空间拟紧。

\begin{theorem}\label{b3:thm:affine-scheme-hom}
设 \(A,B\) 为交换环。取全局截面给出自然双射
\[
\operatorname{Hom}_{\mathrm{Sch}}(\operatorname{Spec}B,\operatorname{Spec}A)
\cong\operatorname{Hom}_{\mathrm{Ring}}(A,B),
\]
并把态射复合对应于环同态的反向复合。
\end{theorem}
\begin{proof}
给定 \(\varphi:A\rightarrow B\)，点映射为 \(\mathfrak q\mapsto\varphi^{-1}(\mathfrak q)\)，连续性已由\cref{b3:prop:spectrum-functor}证明。分式映射 \(a/s\mapsto\varphi(a)/\varphi(s)\) 在相应基本开集上相容，所以给出层映射；茎映射的极大理想原像显然为 \(\varphi^{-1}(\mathfrak q)A_{\varphi^{-1}(\mathfrak q)}\)，故是局部同态。

反之，从概形态射 \((f,f^\#)\) 取全局截面得 \(\varphi:A\rightarrow B\)。对于 \(\mathfrak q\in\operatorname{Spec}B\)，一个 \(a\in A\) 属于 \(f(\mathfrak q)\)，等价于其芽在目标局部环 \(A_{f(\mathfrak q)}\) 的极大理想中。局部同态的极大理想原像条件使这又等价于其像在源局部环 \(B_{\mathfrak q}\) 的极大理想中，即 \(\varphi(a)\in\mathfrak q\)。于是 \(f(\mathfrak q)=\varphi^{-1}(\mathfrak q)\)。在茎上，分母一旦可逆，映射只能是上述分式映射；层态射由所有芽决定，故也唯一。这证明两种构造互逆，复合的断言由拉回次序得到。
\end{proof}

这项对应说明，仿射概形的全局截面环既能恢复整个仿射概形，又能确定它到另一个仿射概形的态射。结构层保留了各局部化，但在仿射情形，这些局部数据仍由同一个全局环控制；一般概形还要保留仿射图之间的黏合数据。

\ProseParagraphBreak
局部性是实质条件。若 \(R\) 为 DVR、\(K\) 为其分式域，把 \(\operatorname{Spec}K\) 的点送到 \(\operatorname{Spec}R\) 的闭点，并在环层上用 \(R\hookrightarrow K\)，虽可得到环空间态射，却不是概形态射。取统一化元 \(\pi\)：它在 \(R\) 的闭点处属于极大理想，拉回到 \(K\) 却成为单位，因而没有保留指定点处的消失条件。等价地，\(K\) 的零极大理想收缩为 \((0)\)，不是 \(R\) 的极大理想。真正由环包含 \(R\hookrightarrow K\) 诱导的概形态射应把唯一点送到一般点 \((0)\)。定理及此条件的辨析见\cite[Chapter II, Proposition 2.3 and Example 2.3.2, pp. 73--74]{Hartshorne1977}。
